\section{Literal self-reference: the completed recurrence}
\label{sec:selfref}
Version 3.6.1.6 adds the executable profile \textbf{SRK-R1}. Its state is a
32-bit word $q_t$ that occurs in its own next-state equation. The old release
provided literal ASA/NA and J--K primitives but its observation adapter used a
constant hold event. This subversion supplies an explicit feedback composition,
editable equation input, CPU/CUDA execution and independent equation replay.

Let $W=\{0,\ldots,2^{32}-1\}$, and let $V,A,N,B,d\in W$ denote the present,
ASA, NA, boundary and external drive words. All complements below are 32-bit
complements. The initial state satisfies $q_0\land\neg V=0$. A trajectory supplies
three Boolean expressions $E_X,E_J,E_K$ and a requested number of steps $T>0$.
The complete transition, in evaluation order, is
\begin{align}
x_t &= E_X(q_t,d),\label{eq:sr-x}\\
a_t &= x_t\land A\land V,\\
n_t &= a_t\land N,\\
h_t &= \pop(n_t\land B),\\
y_t &= \begin{cases}0,&h_t>0,\\ n_t,&h_t=0,\end{cases}\label{eq:sr-asa}\\
J_t &= E_J(q_t,y_t,d),\qquad K_t=E_K(q_t,y_t,d),\\
q_{t+1} &= \big[(J_t\land\neg q_t)\lor(\neg K_t\land q_t)\big]\land V.
\label{eq:sr-jk}
\end{align}
The ASA tuple $(a_t,n_t,h_t,y_t)$ is the retained whole-word rule. The bracketed
part of \eqref{eq:sr-jk} is the retained literal J--K equation. The supplied
expression bindings and final projection to $V$ are explicit SRK-R1 definitions.
For a single active bit they give hold, set, reset and toggle for $(J,K)$ equal
to $(0,0),(1,0),(0,1),(1,1)$ respectively.

\begin{panel}
\textbf{Self-reference is an executable dependency}\par
$q_t\longrightarrow x_t\longrightarrow(a_t,n_t,h_t,y_t)
\longrightarrow(J_t,K_t)\longrightarrow q_{t+1}$, with $q_t$ also read directly
by both J/K expressions and by the final update. Every right-hand side reads the
same old state. Only the final line commits the next state.
\end{panel}
Boundary absorption sets the whole output word $y_t$ to zero. It does not
silently replace the user's J/K expressions with hold: those expressions still
evaluate with that zero word. Thus the equations determine the transition even
in an absorbed step. Numerical GNSS residual classifications do not gate SRK-R1.
\clearpage

\section{Actual equations are input data}
The editable input is \code{examples/self\_reference/equations.json}. The root
declares version \code{3.6.1.6}, profile \code{SRK-R1}, optional default
\code{equations}, and a \code{trajectories} array. Each trajectory supplies an ID,
initial state, drive, four masks and step count; it may override the equations.
For example:
\begin{lstlisting}
{
  "version": "3.6.1.6", "profile": "SRK-R1",
  "equations": {"x": "q | d", "j": "y", "k": "y"},
  "trajectories": [{
    "id": 1, "q0": 0, "drive": 15, "present": 15,
    "asa_mask": 15, "na_mask": 15, "boundary_mask": 0,
    "steps": 6
  }]
}
\end{lstlisting}
The expression vocabulary is \code{\&}, \code{|}, \code{\textasciicircum}, \code{\textasciitilde},
parentheses, the names \code{q}, \code{d}, and, for J/K, \code{y}. The literal
\code{0} means the zero word; \code{1} means the all-ones word. Decimal state and
mask fields remain ordinary unsigned integers. Operator precedence is complement,
AND, XOR, OR; parentheses can specify every grouping explicitly.

The expression trees define an exact finite Boolean algebra. They are parsed
and evaluated as that algebra; they are not passed to a general-purpose evaluator.
Each expression is applied independently to each bit. Whole-word absorption
couples the bits through $h_t$. Arbitrary cross-bit arithmetic, shift operators
and continuous-state formulae would require a separately defined extension.

\subsection{Exact truth-table compilation}
For bits $q_i,y_i,d_i\in\{0,1\}$, the compiler creates
\begin{align}
L_X &= \sum_{u,v\in\{0,1\}} E_X(u,v)\,2^{2u+v},\\
L_J &= \sum_{u,v,w\in\{0,1\}} E_J(u,v,w)\,2^{4u+2v+w},\\
L_K &= \sum_{u,v,w\in\{0,1\}} E_K(u,v,w)\,2^{4u+2v+w}.
\end{align}
The CPU/CUDA kernel reads the corresponding table bit. For the example above,
$(L_X,L_J,L_K)=(14,204,204)$. For the latch with $K=0$, the tuple is
$(14,204,0)$. Every allowed expression has a truth table on these inputs, and
every Boolean function on these inputs is expressible in this vocabulary.
The compilation is exact: each possible input assignment is evaluated, and the
native table read retrieves that assignment's value.
\clearpage

\section{Fixed points, cycles and complete finite execution}
During one trajectory, equations, drive and masks are constant. They define a
deterministic function $F:Q_V\to Q_V$, where
\begin{equation}
Q_V=\{q\in W:q\land\neg V=0\},\qquad |Q_V|=2^{\pop(V)},\qquad
q_{t+1}=F(q_t).
\end{equation}
\begin{proposition}[Closure and eventual periodicity]
Every computed state belongs to $Q_V$. Every infinite orbit of this fixed-input
transition is eventually periodic. A fixed point is a cycle of length one.
\end{proposition}
\begin{proof}
The input requires $q_0\in Q_V$ and the final AND with $V$ places every successor
in $Q_V$. Among $|Q_V|+1$ consecutive states at least two are equal. If
$q_i=q_j$ for $i<j$, determinism gives $q_{i+k}=q_{j+k}$ for every $k\ge0$.
This proves periodicity after the first repeated state; it does not assert
that the period is one.
\end{proof}
The implementation emits every requested step, including those after a repeat.
Independent replay records a fixed point, a cycle, or no repeat observed within
the requested budget. The last case is finite evidence only; it does not imply
divergence or the absence of a later repeat. With 32 active bits the general
state bound is $2^{32}$, so this theorem alone gives no short runtime bound.

\subsection{Worked exact trajectories}
Set $V=A=N=d=15$, $B=0$, $q_0=0$, and $E_X=q\lor d$, $E_J=y$.
For the latch $E_K=0$, $x=a=n=y=15$ and
\begin{equation}
q_{t+1}=y\lor q_t=15,\qquad (q_0,q_1,q_2,\ldots)=(0,15,15,\ldots).
\end{equation}
For toggle $E_K=y$, $J=K=15$ and
\begin{equation}
q_{t+1}=q_t\mathbin{\oplus}15,\qquad
(q_0,q_1,q_2,q_3,\ldots)=(0,15,0,15,\ldots).
\end{equation}
With $B=8$ and $E_J=E_K=y$, $h=1$ absorbs all of $y$ and the state holds.
With the same absorbed word but $E_J=d$, $E_K=0$, the state becomes 15.
These outcomes follow directly from the supplied equations and expose the
difference between output absorption and state reset.

Synchronous evaluation matters. For one active bit with $E_J=\neg q$ and
$E_K=q$, the result is toggle. Both expressions use $q_t$, never a partially
updated state. State projection preserves inactive bits as zero, including when
an input expression uses the all-ones constant.
\clearpage

\section{Native execution and independent equation replay}
The shared \code{include/self\_reference.hpp} implements the truth-table and
literal transition operations. The CPU walks each trajectory in step order.
CUDA assigns one thread to each independently seeded trajectory, then executes
its recurrence sequentially within that thread. No ordering between independent
trajectories is needed; different trajectories never share a mutable state.

The compilation interface is a decimal CSV with columns
\begin{lstlisting}
trajectory_id,q0,drive,present,asa_mask,na_mask,boundary_mask,
x_lut,j_lut,k_lut,steps
\end{lstlisting}
The header is physically one line. Every trace row retains
\begin{lstlisting}
trajectory_id,step,before,drive,x,asa,na,hits,output,j,k,after
\end{lstlisting}
The native run also records backend, device, timing and trace lengths. The
Python reference evaluates the original expression trees on whole integer words;
it does not use the compiled truth tables as its transition oracle. It compares
every trace field, row count and trajectory ordering, and records equation/input
provenance. This cross-check can detect table compilation, transition and output
corruption; it is finite execution evidence rather than a machine-checked proof.

\subsection{Reproduce the equation-input run}
From the package directory, build CPU and run:
\begin{lstlisting}
cmake -S . -B build_cpu
cmake --build build_cpu --config Release --parallel 2
python tools/self_reference.py \
  --input examples/self_reference/equations.json \
  --out results/selfref_local \
  --binary build_cpu/Release/self_reference.exe --backend cpu
\end{lstlisting}
For single-configuration systems the binary is \code{build\_cpu/self\_reference}.
For CUDA, enable \code{SATNAV\_ENABLE\_CUDA=ON}, choose the GPU binary and use
\code{--backend cuda}. Each run uses a new output directory. The JSON expressions
are editable without recompiling C++ or CUDA; the compiler regenerates their tables.

\subsection{Relation to the existing SATNAV and UGTS equations}
SATNAV-R1 retains its corrected-code equations, QR solve, OTAN2 source ratio,
both 64-bit codecs and full coordinate/time records. The earlier observation adapter's
hold event is retained by the explicit legacy option. SRK-R1 is the retained formal word
state-feedback computation and runs without a GNSS residual admission condition.
The new bindings are identified as definitions of this release. CGK-R1 now provides explicit new circle-plus, encoder and physical-matrix
definitions and connects the observation path, as detailed in Section~\ref{sec:coupled}.
The unavailable earlier M2 definitions are not claimed to have been reconstructed.
\clearpage

\section{SRK-R1 baseline validation, version 3.6.1.6}
All entries here concern the executable 3.6.1.6 source on the current Windows
machine. The 3.6.1.5 preparation record is preserved separately under
\code{results\_3\_6\_1\_5/}. Finite comparisons check the requested trajectories;
the closure and eventual-periodicity argument is given separately above.

\begin{center}
\begin{tabularx}{\textwidth}{@{}P{.53\textwidth}Y@{}}\toprule
\textbf{Executed check} & \textbf{Result}\\\midrule
CPU CMake build and CTest & Passed, 2 test programs\\
Retained SATNAV C++ checks & 105,559 passed\\
New literal feedback C++ checks & 5,060,172 passed\\
Python test methods & 31 passed\\
Native command-line scenarios & 25 feedback + 11 SATNAV passed\\
Editable six-trajectory example & 34 transitions, 408 exact field comparisons per backend\\
257-trajectory stress input & 4,225 transitions, 50,700 exact field comparisons per backend\\
CPU and CUDA equation replay & Both passed sample and stress inputs\\
Feedback GPU Compute Sanitizer & Zero memory errors; sanitized trace replay passed\\
Retained SATNAV independent reference & 256 epochs, 4,828 comparisons passed\\
SATNAV CUDA texture/global paths & Both passed independent comparisons and memcheck\\
CPU sanitizers in this Windows run & Not run; earlier Linux evidence archived\\\bottomrule
\end{tabularx}
\end{center}

\subsection{Recorded toolchain and device}
MSVC 19.44.35221.0, Visual Studio 2022 BuildTools; NVIDIA CUDA compiler 12.8.61;
NVIDIA GeForce RTX 5070 Ti Laptop GPU, compute capability 12.0, driver 591.59.
The source targets \code{sm\_120} and \code{compute\_120}. Device, runtime,
allocation sizes and kernel-only times are retained in each native run record.
The short external build directories used on this machine avoid Windows compiler
scratch-path limits caused by the long workspace path.

The stress input uses seed 3616, varying 1--32 steps per trajectory, full 32-bit
operands and 257 trajectories so that the 128-thread CUDA blocks include a partial
final block. Direct AST replay checks all 12 trace columns. The Python suite also
rejects deliberate mutations of every trace column, missing/extra rows and reordered
rows. The C++ checks include independent bit oracles and exhaustive small domains.

\subsection{Evidence and delivered implementation}
\code{source/baseline\_3\_6\_1\_6\_validation.json} indexes the archived baseline evidence. The exact equation
JSON, compiled CSV, complete native/reference traces and their hashes are retained.
The baseline delivered Windows CPU and CUDA executables under \code{bin/windows/}; CUDA
execution uses the installed NVIDIA runtime and driver. No measured speedup is
inferred from these small fixtures. The PDF is built with Tectonic 0.17.0 and
visually checked in that release. Source and delivered artifacts were listed in
the regenerated SHA-256 manifest.
\clearpage

